\documentclass[12pt]{article}
\usepackage{amsmath,amssymb,amsthm}
\usepackage[margin=1in]{geometry}

\title{Project Euler Problem 866: Tidying Up B}
\author{Solution}
\date{}

\begin{document}
\maketitle

\section{Problem Statement}
A small child has a ``number caterpillar'' consisting of $N$ jigsaw pieces numbered $1$ to $N$. The father picks up pieces at random and places them in their correct positions. When a segment of length $k$ is formed, the $k$-th hexagonal number $H_k = k(2k-1)$ is recorded. The product of all recorded numbers is computed.

Find the expected value of this product for $N = 100$, modulo $987654319$.

\section{Mathematical Analysis}

\subsection{Hexagonal Numbers}
The $k$-th hexagonal number is:
\[
H_k = k(2k - 1)
\]

\subsection{Segment Dynamics}
When piece at position $i$ is placed (at some random time), the following cases arise:
\begin{itemize}
    \item \textbf{Isolated}: No neighbors filled. New segment of length 1. Record $H_1 = 1$.
    \item \textbf{Extend}: One neighbor filled, extending segment from length $m$ to $m+1$. Record $H_{m+1}$.
    \item \textbf{Merge}: Both neighbors filled, merging segments of lengths $m_1, m_2$ into $m_1 + m_2 + 1$. Record $H_{m_1+m_2+1}$.
\end{itemize}

\subsection{Expected Product via DP}
We define a DP over intervals. Consider building the caterpillar by inserting pieces into $[1, N]$. The key insight is that the expected product can be decomposed by considering which piece is placed \emph{last} in each contiguous interval.

For an interval $[a, b]$ of length $L = b - a + 1$, define $E(L)$ as the expected product of hexagonal numbers generated when filling this interval. The piece placed last in the interval contributes $H_L$ to the product, and the two sub-intervals to its left and right contribute independently.

\[
E(L) = H_L \cdot \frac{1}{L} \sum_{i=1}^{L} \binom{L-1}{i-1} \cdot \frac{(i-1)!(L-i)!}{(L-1)!} \cdot E(i-1) \cdot E(L-i)
\]

Simplifying:
\[
E(L) = \frac{H_L}{L} \sum_{i=1}^{L} E(i-1) \cdot E(L-i)
\]

with $E(0) = 1$.

\subsection{Modular Arithmetic}
All computations are performed modulo $p = 987654319$, which is prime. We use Fermat's little theorem for modular inverses: $a^{-1} \equiv a^{p-2} \pmod{p}$.

\section{Algorithm}
\begin{enumerate}
    \item Compute $E(L)$ for $L = 0, 1, \ldots, 100$ using the recurrence above.
    \item The answer is $E(100) \bmod 987654319$.
\end{enumerate}

\section{Complexity}
\begin{itemize}
    \item \textbf{Time}: $O(N^2)$ for the DP.
    \item \textbf{Space}: $O(N)$.
\end{itemize}

\section{Answer}

\[
\boxed{492401720}
\]

\end{document}
